%
% ---------- header -----------------------------------------------------------
%
% project       kaneton
%
% license       kaneton
%
% file          /home/mycure/kane...view/book/architecture/mips64/mips64.tex
%
% created       enguerrand raymond   [wed oct 01 19:07:46 2008]
% updated       julien quintard   [mon jan 31 14:28:33 2011]
%

%
% ---------- setup ------------------------------------------------------------
%

%
% path
%

\def\path{../../..}

%
% template
%

%
% ---------- header -----------------------------------------------------------
%
% project       kaneton
%
% license       kaneton
%
% file          /home/mycure/kaneton/view/template/book.tex
%
% created       julien quintard   [wed may 16 18:16:50 2007]
% updated       julien quintard   [sun jun  3 12:00:37 2007]
%

%
% class
%

\documentclass[10pt,a4wide]{book}

%
% packages
%

\usepackage[english]{babel}
\usepackage[T1]{fontenc}
\usepackage{a4wide}
\usepackage{fancyheadings}
\usepackage{multicol}
\usepackage{indentfirst}
\usepackage{graphicx}
\usepackage{color}
\usepackage{xcolor}
\usepackage{verbatim}

\usepackage{aeguill}

\usepackage[Lenny]{\path/package/fncychap}

\pagestyle{fancy}

%
% -
%

\renewcommand{\-}{\vspace{\parskip}}

%
% toc
%

\newcommand{\toc}
  {
    \tableofcontents
    \setlength{\footrulewidth}{0.3pt}
    \setlength{\parindent}{0.3cm}
    \setlength{\parskip}{2ex plus 0.5ex minus 0.2ex}
  }

%
% logos
%

\newcommand{\logos}
  {
    \begin{center}
      \includegraphics[scale=0.8]{\path/logo/kaneton.pdf}
    \end{center}
  }

%
% colors
%

\definecolor{functioncolor}{rgb}{0.40,0.00,0.00}
\definecolor{commandcolor}{rgb}{0.00,0.00,0.40}
\definecolor{verbatimcolor}{rgb}{0.00,0.40,0.00}
\definecolor{noticecolor}{rgb}{0.87,0.84,0.02}

%
% function
%

\newcommand\function[3]{
  \begin{tabular}{p{0.2cm}p{13.8cm}}
  & {\color{functioncolor}\textbf{#1}}#2
  \end{tabular}

  \begin{tabular}{p{1cm}p{13cm}}
  & #3
  \end{tabular}}

%
% align
%

\newcommand\align[1]{
  \\ & \hspace{#1}}

%
% argument
%

\newcommand\argument[1]{\textit{#1}}

%
% command
%

\newcommand\command[2]{
  \begin{tabular}{p{0.2cm}p{13.8cm}}
  & {\color{commandcolor}\textbf{#1}}
  \end{tabular}

  \begin{tabular}{p{1cm}p{13cm}}
  & #2
  \end{tabular}}

%
% notice
%

\newcommand\notice[1]{
  {\color{noticecolor}\textbf{Notice}}

  \begin{tabular}{p{0.2cm}p{13.8cm}}
  & #1
  \end{tabular}}

%
% subsubsubsection
%

\newcommand\subsubsubsection[1]{\textbf{#1}}

%
% example
%

\newcommand\example[1]{
  \textit{Example:}

  \begin{tabular}{p{0.2cm}p{13.8cm}}
  & \textit{#1}
  \end{tabular}}

%
% warning XXX
%

\renewcommand{\familydefault}{\sfdefault}

%
% verbatim stuff
%

\makeatletter

\renewcommand{\verbatim@font}
  {\ttfamily\footnotesize\selectfont}

\def\verbatim@processline{
  \hskip15ex{\color{verbatimcolor}\the\verbatim@line}\par
}

\makeatother

%
% header
%

\rhead{}
\rfoot{\scriptsize{The kaneton microkernel project}}

\date{\scriptsize{\today}}


%
% header
%

\lhead{\scriptsize{The kaneton microkernel :: mips64}}

%
% title
%

\title{The kaneton microkernel :: mips64 \\
       \version
       \logo}

%
% document
%

\begin{document}

%
% title
%

\maketitle

%
% --------- text --------------------------------------------------------------
%

This document describes the kaneton microkernel reference project over
MIPS64 architecture on qemu platform (for educational purpose).

\-

This document should be used by every student willing implement the
kaneton microkernel as people looking for more details on the kaneton
microkernel implementation on MIPS64 architecture. Several microprocessor 
kaneton implementations will be described inside this document.

\-

All the kaneton documents are available on
the official website
  \footnote{\location{http://kaneton.opaak.org}}.

\-

\textbf{Prerequisites}: this document details the kaneton implementation
for MIPS64 architecture on Qemu platform (for educational purpose). A
knowledge of this architecture's basics is strongly recommended as
this document is not intended to be a MIPS64 lesson (see
\textit{Bibliography} for references).

%
% table of contents
%

\tableofcontents

%
% identation
%

\indentation{}

%
% chapters
%

%%
%% licence       kaneton licence
%%
%% project       kaneton
%%
%% file          /home/mycure/kaneton/view/papers/kaneton/overview.tex
%%
%% created       matthieu bucchianeri   [mon jan 30 17:09:45 2006]
%% updated       julien quintard   [thu mar  2 13:12:22 2006]
%%

%
% overview
%

\chapter{Overview}

XXX ce chapitre va vous aider a reconnaitre les fonctionnalites principale
XXX d'un kernel dans kaneton.

The kaneton microkernel is only the core of an operating system.
Main tasks like hardware drivers or user services are implemented as
\textbf{servers}. So the microkernel only has a few functionalities to
provide:

\begin{itemize}
  \item
    Memory management.
  \item
    Process management.
  \item
    Communication.
  \item
    Events.
\end{itemize}

In this chapter we will describe briefly these tasks and all the
associated managers.

%
% memory management
%

\section{Memory Management}

Handling the memory -- from virtual address space to physical
addressing -- is done by three major managers, the \textbf{as},
\textbf{segment} and \textbf{region} managers.

%
% as
%

\subsection{as}

The address space manager just manages the different address spaces
used by the kaneton tasks.

In kaneton, we call an \textbf{as - address space} a list of memory
locations referenced by a task. Each task has its own address space.

%
% segment
%

\subsection{segment}

The segment manager just manages the segments reserved by
the different kaneton entities including the kernel, the drivers etc..

In kaneton terms a \textbf{segment} is a contiguous area of reserved
physical memory.

%
% region
%

\subsection{region}

The region manager keeps track of regions used to map segments for
each address space reserved on the system.

In kaneton, a \textbf{region} is contiguous area of virtual memory
mapping a segment's part.

%
% process management
%

\section{Process Management}

XXX

%
% communication
%

\section{Communication}

XXX

%
% events
%

\section{Events}

XXX

%
% ---------- header -----------------------------------------------------------
%
% project       kaneton
%
% license       kaneton
%
% file          /home/enguerrand/...ew/book/architecture/MIPS64/platform.tex
%
% created       enguerrand raymond   [mon oct 13 11:31:56 2008]
% updated          [fri mar 27 15:18:49 2009]
%

%
% platform
%

\chapter{MIPS64 platforms}

This chapter describes the different platform for which kaneton has been developed. These platforms are described following you need to understand the current document.


\subsection{qemu-mips platform}

\subsubsection{Description}

In 2008 We did not have MIPS machine enough in EPITA laboratories. We have decided to develop the first MIPS64 kaneton
port on MIPS64 emulator only. The student must already develop on ia32 architecture and they must execute their kernel
on real machine. So, it is not necessary to have a real MIPS 64 bits machine.

The qemu-mips platform is the simplest emulated platform to develop kernel for MIPS 64 bits architecture. So, we have decided to develop firstly on this platform. It is composed of the following devices :

\begin{itemize}
  \item
    \textbf{A big range of MIPS64 microprocessor}
  \item
    \textbf{IBM PC style serial port}
  \item
    \textbf{IBM PC style IDE disk}
  \item
    \textbf{NE2000 network card}
\end{itemize}
 
\subsubsection{Emulation launching}

Qemu-mips does not have any firmware (on MIPS platform the bios is called a firmware).

We use the qemu-mips platform with the R4000 microprocessor. The command line to launch him it the following :

\begin{verbatim}
  qemu-system-mips64el [-L path ot bios] -fda kaneton.img -M mips -cpu R4000
\end{verbatim}

There is no mips firmware (or bios for qemu) today but we will launch the kernel as it. It will load at the address \code{0xFFFFFFFFBFC00000}. To do that, the kernel must be copied to \code{mips\_bios.bin} and the \code{L} option must indicate the path to this file. We will develop a firmware later.
The bootloader will be executed like a firmware.

\subsubsection{MIPS R4000 microprocessor}

The first 64 bits MIPS microprocessor was the R4000, it released in 1991 but MIPS could not produce it themselves.
In 1992 SGI bought the MIPS company that becam MIPS technologies. Thus the microprocessor could be produce for SGI platforms. The MIPS R4000 has been designed in three configurations for low-cost system (R4000PC), high performance uniprocessor system (R4000SC) and for high multiprocessor platform (R4000MC).

Whe have chosen this processor because it is the simplest MIPS 64 bits microprocessor and it has all functionalities that we need for the pedagogical purposes.

%%
%% licence       kaneton licence
%%
%% project       kaneton
%%
%% file          /home/mycure/kaneton/view/books/kaneton/development-environment.tex
%%
%% created       matthieu bucchianeri   [mon jan 30 17:31:26 2006]
%% updated       julien quintard   [sun jul 30 23:01:51 2006]
%%

%
% development environment
%

\chapter{Development Environment}

In this chapter we will study the kaneton development environment
to understand how it works and how to extend it with new users and
machines profiles.

\newpage

%
% text
%

Over the years, the kaneton microkernel evolved, starting with a
very simple introduction to low-level programming and finally
to a complete microkernel development.

Going always further implies many modifications. Indeed, the students
cannot take time to develop their own development environment because
developing a development environment is a whole project by itself.

We wanted to lead students to a complete microkernel development and then
to introduce the parallel and distributed programming. This would not
be possible if they had to develop an entire development environment.

So kaneton people decided to give the students a complete development
environment.

The kaneton development environment is composed of makefiles, shell
scripts and configuration files and together provide multi-users cooperation,
with many operating systems on many architectures.

A development environment describes binaries, compiling options, linking
options etc.. to use for the kaneton microkernel project.

The whole kaneton development environment needs exactly two
fundamental tools to work. The first one is GNU make, used to build
powerful makefiles, and the second one is GNU bash, used to write
shell scripts.

If an operating system has these two tools, then kaneton can be developed
on it.

%
% profiles
%

\section{Profiles}

The kaneton development environment permit different users, with different
development parameters and configuration files, running different operating
systems on different architectures, to develop together on the same
microkernel project.

To offer these functionalities, the kaneton development environment is
composed of profiles, machine profiles which describe operating systems
behaviours and user profiles which describe users configurations.

From these profiles, the kaneton development environment will generate
development environment files containing the whole kaneton environment
stuff including variables, shell functions, make functions etc..

The three kaneton development environment files are located in the
\textit{env/} directory.

The first one contains static kaneton development environment information
and is called \textit{.env.conf}. This file contains absolute paths
to different kaneton directories and kaneton general information.

The two others called \textit{.env.mk} and \textit{.env.sh} will be generated
by the kaneton development environment system. These files contain
everything necessary to kaneton makefiles and kaneton shell scripts
including variables and functions.

The kaneton development environment system uses shell environment variables
the user has to set before using the system:

\begin{itemize}
  \item
    \textit{\$KANETON\_USER} holds the user profile name.

    This shell environment variable will be used by the kaneton development
    environment system to retrieve the user profile.
  \item
    \textit{\$KANETON\_MACHINE} holds the machine profile name.

    This value will be used to retrieve the machine profile.
  \item
    \textit{\$KANETON\_SHELL} holds the GNU bash binary path.

    This path will be used to launch the very first shell script used to
    generate the current development environment.
\end{itemize}

%
% machines
%

\subsection{Machines}

kaneton people wanted a powerful system to provide users an easy way
to use kaneton on any operating system and any microprocessor architecture.

Indeed, many students and developers use Apple's MacOS running on PowerPC
microprocessors while other use Solaris running on UltraSparc microprocessors
and finally the others generally use Linux or BSD operating systems on
Intel microprocessors.

kaneton people wanted all the developers to contribute to the same
kaneton implementation without difficulties.

To make it possible, the kaneton microkernel had to compile on
every operating systems. To do so, machine profiles were introduced to
describe operating systems behaviour.

A machine profile is precisly a couple operating system, running
microprocessor architecture and kaneton target microprocessor architecture.

The naming scheme used for machine profiles is the operating system name,
a dash, the running on microprocessor, a dot and the target microprocessor
architecture for which kaneton is built.

Note that in the case of a same running on microprocessor and target
microprocessor, only the first naming part is necessary.

For example, a correct machine profile name would be
\textit{macos-ppc.ia32} describing the machine profile composed of the MacOS
operating system running on PowerPC microprocessor architecture and generating
kaneton binary object files for the Intel Architecture 32-bit. Another correct
machine profile name would be \textit{freebsd} describing the behaviour of
the FreeBSD operating system running on any microprocessor and generating
kaneton binary object files for the same microprocessor's architecture.

The machine profiles are located in the \textit{env/machines/} directory.
A subdirectory is used for every machine profiles.

Notice that the kaneton development environment will use the shell
environement variable \textit{\$KANETON\_MACHINE} to locate the
machine profile to use.

Then the current machine profile's files are always located in the
\textit{env/machines/\$KANETON\_MACHINE/} directory.

A machine profile describes the operating system behaviour for very
generic operations like compiling a file, assembling a file, building
a link between two files, removing a file, changing the current directory,
displaying a message etc..

As seen earlier, the kaneton development environment describes two
behaviours, one for the makefiles, and another for the shell scripts.

The files located in each machine profile are listed below. Note that
machine profiles may add any necessary file in their directory.

\begin{verbatim}
clean.sh
critical.sh
init.sh
machine.conf
machine.mk
machine.sh
\end{verbatim}

The files \textit{critical.sh}, \textit{init.sh} and \textit{clean.sh} are
used to install and uninstall the development environment for this
machine profile.

More precisly, the \textit{critical.sh} shell script must generate
the development environment fundamental files \textit{.env.mk} and
\textit{.env.sh} located in the \textit{env/} directory.

Then, the \textit{init.sh} is called to install machine specific stuff
like links to machine dependent directories. This shell script should also
check for the presence of binaries used by the machine profile.

Finally, the \textit{clean.sh} is called to uninstall the machine
specific stuff.

The \textit{machine.conf} file describes the machine profile configuration
including general variables, general flags etc.. for the machine profile.

The file \textit{machine.mk} describes the machine behaviour for
make operations. This file must implement the whole kaneton make interface.

The file \textit{machine.sh} describes the machine behaviour for
shell operations. This file must implement the whole kaneton shell
interface.

Note that it might exist partial behaviours while the user willing use
this machine profile must implement the other part.

Then, with all these files, an operating system can be described to
perform very basic make and shell operations generating binary object
files for the current microprocessor architecture or for another specific
one.

%
% users
%

\subsection{Users}

Each kaneton developer has his own profile located in the kaneton
hierarchy \textit{env/users/\$KANETON\_USER/}.

Note that this shell environment variable must be composed of a
first name, a dot and a last name; using nickname is not allowed.

For example, a correct user name would be \textit{julien.quintard} and
a bad one would be \textit{ChIcHe}.

The user profile directory contains configuration files used to parameter
the developement environment and the kaneton microkernel. Since every
users has his own profile, the current developer will use his own user
profile configuration and will not interfer with other developers.

The files located in each user profile directory are listed below.

\begin{verbatim}
  conf.c
  conf.h
  kaneton.conf
  modules.conf
  user.conf
  user.mk
  user.sh
\end{verbatim}

The file \textit{conf.c} is not used in the current development environment
implementation.

The file \textit{conf.h} parameterizes the kaneton microkernel specifying
for example the algorithm to use in the segment and region managers,
but is also used to activate for example the debug modes etc..

The file \textit{kaneton.conf} is the runtime kaneton microkernel
configuration file. This file describes the kaneton servers to launch
at the operating system startup. Moreover this file is used to parameter
the kaneton microkernel and servers in a dynamic way.

The file \textit{modules.conf} contains the module files to pass to
kaneton. Note that the first module in this file will be considered as
the kaneton bootloader and the second one as the kaneton core.

The file \textit{user.conf} is the most important since it is used
to configure the whole kaneton development environment. This file uses
the GNU make syntax. Via this file, the developer can parameter the compiling
flags, the boot device modes, the display options, the microkernel
files to use etc..

The \textit{user.mk} makefile is used to redefine or complete the machine
makefile behaviour.

Indeed, a developer could want not to use the general machine behaviour but
instead, prefer redefine parts of it. Then the user can override parts of
the machine profile's makefile behaviour using this file to redefine makefile
functions.

The \textit{user.sh} shell file is used, as the \textit{user.mk}, to
redefine shell functions.

With all these files a kaneton developer can parameterize the whole
kaneton project.

%
% interfaces
%

\section{Interfaces}

%
% make
%

\subsection{Make}

In this section we will detail the make interface that every machine
profile must implement.

The reader should look closer to the machine profiles already implemented.

The current machine reference profile is the \textit{linux} one.

\function{print}{(\argument{color},
                  \argument{text},
                  \argument{options})}
         {
	   This function display the \argument{text} in the specified
	   \argument{color}.

	   The option \textit{--no-newline} can be used not to output
	   the trailing newline.
	 }

\function{change-directory}{(\argument{directory},
                             \argument{options})}
         {
	   This function changes the current working directory.
	 }

\function{launch-shell}{(\argument{shell script},
                         \argument{arguments},
                         \argument{options})}
         {
	   This function launches a new shell script with its arguments.
	 }

\function{launch-python}{(\argument{python script},
                          \argument{arguments},
                          \argument{options})}
         {
	   This function launches a new python script with its arguments.
	 }

\function{launch-perl}{(\argument{shell script},
                        \argument{arguments},
                        \argument{options})}
         {
	   This function launches a new perl script with its arguments.
	 }

\function{launch-make}{(\argument{make file},
                        \argument{directory list},
                        \argument{arguments},
                        \argument{options})}
         {
	   This function launches the make command in each directory of
	   the list.
	 }

\function{launch}{(\argument{file},
                   \argument{...})}
         {
	   This function is a wrapper over the functions to launch new
	   shell, python, perl and make scripts.
	 }

\function{preprocess}{(\argument{preprocessed file},
                       \argument{c file},
                       \argument{options})}
         {
	   This function launches the C preprocessor.
	 }

\function{compile-c}{(\argument{object file},
                      \argument{c file},
                      \argument{options})}
         {
	   This function compile a C file generating an object file.
	 }

\function{lex-l}{(\argument{c file},
                  \argument{lex file},
                  \argument{options})}
         {
	   This function generates a C source file from a lex file.
	 }

\function{assemble-S}{(\argument{object file},
                       \argument{S file},
                       \argument{options})}
         {
	   This function assemble a S file.
	 }

\function{assemble-asm}{(\argument{object file},
                         \argument{asm file},
                         \argument{options})}
         {
	   This function assemble an asm file.

	   The option \textit{--output-format [format]} can be used to
	   specify the output format \textit{elf}, \textit{bin} etc..
	 }

\function{static-library}{(\argument{static library file name},
                           \argument{object files},
                           \argument{options})}
         {
	   This function builds a static library from object files.
	 }

\function{dynamic-library}{(\argument{dynamic library file name},
                            \argument{object files and/or libraries},
                            \argument{options})}
         {
	   This function builds a dynamic library from object files and/or
	   libraries.
	 }

\function{executable}{(\argument{executable file name},
                       \argument{object files and/or libraries},
                       \argument{options})}
         {
	   This function builds a executable from object files and/or
	   libraries.

	   The option \textit{--link-script [file]} specifies the
	   linker script file to use when building the executable object.

	   The option \textit{--no-standard-include} tells the function
	   not to use the operating system standard includes.

	   The option \textit{--no-standard-library} tells the function
	   not to use the operating system standard library.
	 }

\function{archive}{(\argument{archive file name},
                    \argument{object files},
                    \argument{options})}
         {
	   This function builds an archive of object files from multiple
	   object files.
	 }

\function{remove}{(\argument{files},
                   \argument{options})}
         {
	   This function removes the files in the list.
	 }

\function{purge}{()}
         {
	   This function just cleans the current working directory from
	   unecessary files.
	 }

\function{prototypes}{(\argument{files},
                       \argument{options})}
         {
	   This function will generate the prototypes in relation with
	   the files in the list.
	 }

\function{dependencies}{(\argument{files},
                         \argument{output},
                         \argument{options})}
         {
	   This function will generate the dependencies for the
	   \argument{files} and build a makefile dependency
	   \argument{output}.
	 }

\function{version}{(\argument{file})}
         {
	   This function generates a version file from the operating
	   system's informations: user, host, date etc..
	 }

\function{link}{(\argument{link},
                 \argument{file},
                 \argument{options})}
         {
	   This function creates a link \argument{link} for the file
	   \argument{file}.
	 }

\function{compile-tex}{(\argument{file},
                        \argument{options})}
         {
	   This function compiles the file \argument{file}.tex and
	   will generate a readable document.
	 }

\function{paper}{(\argument{file},
                  \argument{options})}
         {
	   This function builds a paper by calling the
	   \textbf{compile-tex()} function.
	 }

\function{lecture}{(\argument{file},
                    \argument{options})}
         {
	   This function builds a lecture document by calling the
	   \textbf{compile-tex()} function.
	 }

\function{subject}{(\argument{file},
                    \argument{options})}
         {
	   This function builds a subject by calling the
	   \textbf{compile-tex()} function.
	 }

\function{correction}{(\argument{file},
                       \argument{options})}
         {
	   This function builds a correction document by calling the
	   \textbf{compile-tex()} function.
	 }

\function{view}{(\argument{file},
                 \argument{options})}
         {
	   This function launches a viewer for the readable document
	   generated by the function \textbf{compile-tex()}.
	 }

\function{contents}{(\argument{file},
                     \argument{options})}
         {
	   This function returns the contents of the file \argument{file}.
	 }

%
% shell
%

\subsection{Shell}

In this section we will detail the shell interface that every machine
profile must implement.

The reader should look closer to the machine profiles already implemented.

\function{print}{(\argument{color},
                  \argument{text},
                  \argument{options})}
         {
	   This function outputs some text to the screen using a specific
	   color \textit{red}, \textit{green}, \textit{blue} etc..

	   The option \textit{--no-newline} can be used not to output
	   the trailing newline.

	   The option \textit{--flickering} tells the function to make
	   the text flicker.
	 }

\function{contents}{(\argument{file},
                     \argument{options})}
         {
	   This function returns the contents of the file \argument{file}.
	 }

\function{temporary}{(\argument{options})}
         {
	   This function creates a temporary file system object.

	   The option \textit{--file} tells the function to create a
	   temporary file.

	   The option \textit{--directory} tells the function to create
	   a temporary directory.
	 }

\function{working-directory}{(\argument{options})}
         {
	   This function returns the path of the current working directory.
	 }

\function{locate}{(\argument{name},
                   \argument{options})}
         {
	   This function tries to locate the program \argument{name}
	   on the system.
	 }

\function{substitute}{(\argument{pattern},
                       \argument{string},
                       \argument{options})}
         {
	   This function substitutes any string of the standard input
	   matching the pattern \argument{pattern} by the string
	   \argument{string}.

	   The option \textit{--everywhere} tells the function to
	   maximize the substitutions on the same line.
	 }

\function{display}{(\argument{message},
                    \argument{header})}
         {
	   This function displays the text \argument{message} in a
	   very specific manner based on the argument \argument{header}:
	   \textit{+}, \textit{!}, \textit{?} etc..
	 }

\function{wait-key}{(\argument{options})}
         {
	   This function just waits for the user to press a key.
	 }

\function{launch-shell}{(\argument{file},
                         \argument{arguments},
                         \argument{options})}
         {
	   This function launches a new shell script \argument{file} with
	   the arguments \argument{arguments}.
	 }

\function{launch-python}{(\argument{file},
                          \argument{arguments},
                          \argument{options})}
         {
	   This function launches the python file \argument{file} with
	   the arguments \argument{arguments}.
	 }

\function{launch-perl}{(\argument{file},
                        \argument{arguments},
                        \argument{options})}
         {
	   This function launches the perl file \argument{file} with
	   the arguments \argument{arguments}.
	 }

\function{launch-make}{(\argument{file},
                        \argument{directories},
                        \argument{arguments},
                        \argument{options})}
         {
	   This function launches the make file \argument{file} in each
	   directory of \argument{directories} with the arguments
	   \argument{arguments}.
	 }

\function{launch-program}{(\argument{file},
                           \argument{arguments},
                           \argument{options})}
         {
	   This function launches the program file \argument{file} with
	   the arguments \argument{arguments}.
	 }

\function{launch}{(\argument{file},
                   \argument{...})}
         {
	   This function is a wrapper over the previously described
	   launch functions.
	 }

\function{copy}{(\argument{source},
                 \argument{destination},
                 \argument{options})}
         {
	   This function copies the file \argument{source} to
	   \argument{destination}.
	 }

\function{link}{(\argument{source},
                 \argument{destination},
                 \argument{options})}
         {
	   This function builds a link between the file \argument{source}
	   and the file \argument{destination}.
	 }

\function{remove}{(\argument{file},
                   \argument{options})}
         {
	   This function removes the file \argument{file}.
	 }

\function{list}{(\argument{directory},
                 \argument{options})}
         {
	   This function lists the file system objects contains in the
	   directory \argument{directory}.

	   The option \textit{--file} tells the function to list the
	   files only.

	   The option \textit{--directory} tells the function to list
	   the directories only.
	 }

\function{change-directory}{(\argument{directory},
                             \argument{options})}
         {
	   This function changes the current working directory to
	   \argument{directory}.
	 }

\function{preprocess}{(\argument{file},
                       \argument{includes},
                       \argument{options})}
         {
	   This function performs a preprocessing task on the file
	   \argument{file}.

	   The argument \argument{includes} is used to provide additional
	   includes.

	   The option \textit{--no-markers} inhibits the generation
	   of linemarkers in the output.
	 }

\function{search}{(\argument{path},
                   \argument{pattern},
                   \argument{options})}
         {
	   This function looks for a file given the pattern
	   \argument{pattern} from the directory \argument{path}.

	   The option \textit{--file} is used to retrieve only files.

	   The option \textit{--directory} is used to retrieve only
	   directories.
	 }

\function{pack}{(\argument{directory},
                 \argument{file},
                 \argument{options})}
         {
	   This function makes an archive \argument{file} of the
	   directory \argument{directory}.
	 }

\function{unpack}{(\argument{directory},
                   \argument{file},
                   \argument{options})}
         {
	   This function extractes the archive \argument{file} into the
	   directory \argument{directory}.
	 }

\function{make-directory}{(\argument{directory},
                           \argument{options})}
         {
	   This function builds a new directory named
	   \argument{directory}.
	 }

\function{make-directory}{(\argument{directory},
                           \argument{options})}
         {
	   This function builds a new directory named
	   \argument{directory}.
	 }

\function{cut}{(\argument{head},
                \argument{tail})}
         {
	   This function tries to locate text between the two patterns
	   \argument{head} and \argument{tail} from the standard input.

	   The option \textit{--delete} tells the function to delete
	   this text and display the rest.

	   The option \textit{--print} tells the function to print
	   this text.
	 }

\function{load}{(\argument{source},
                 \argument{device},
                 \argument{path},
                 \argument{options})}
         {
	   This function loads the file \argument{source} on the
	   specificed device \argument{device}, more precisly
	   at the location \argument{path}.

	   The option \textit{--image} specifies that the argument
	   \argument{device} represents a device image.

	   The option \textit{--device} specifies that the argument
	   \argument{device} represents a device.
	 }

\function{stamp}{(\argument{format},
                  \argument{options})}
         {
	   This function returns a date following the given format
	   \argument{format}.
	 }

\function{stamp}{(\argument{format},
                  \argument{options})}
         {
	   This function returns a date following the given format
	   \argument{format}.
	 }

\function{record}{(\argument{log},
                   \argument{time},
                   \argument{options})}
         {
	   This function starts recording a session and outputs
	   the text into the file \argument{log} while the timings
	   are output in the file \argument{time}.
	 }

\function{play}{(\argument{log},
                 \argument{time},
                 \argument{options})}
         {
	   This function plays a previously recorded session where
	   the files \argument{log} and \argument{time} hold the
	   text and the timings.
	 }

%
% actions
%

\section{Actions}

Almost every different actions can be triggered from the kaneton
root directory.

These actions are listed below.

\command{make init}
        {
	  This action will initialize the kaneton development environment.

	  \example{\$ make init}
	}

\command{make clean}
	{
	  This action will clean the kaneton development environment.

	  \example{\$ make clean}
	}

\command{make kaneton}
	{
	  This action will build the kaneton microkernel.

	  \example{\$ make kaneton}
	}

\command{make clear}
	{
	  This action will removing every compiled files.

	  \example{\$ make clear}
	}

\command{make purge}
	{
	  This action will clean directories from unecessary files.

	  \example{\$ make purge}
	}

\command{make dep}
	{
	  This action will generate makefiles dependencies.

	  \example{\$ make dep}
	}

\command{make proto}
	{
	  This action will generate C prototypes.

	  \example{\$ make proto}
	}

\command{make check}
	{
	  This action will be used to run the test suite used to validate
	  the kaneton whole behaviour.

	  \example{\$ make check}
	}

\command{make cheat}
	{
	  This action will launch the cheat tests on students kaneton
	  implementations.

	  \example{\$ make cheat}

	  \example{\$ make cheat-2006.k3}
	}

\command{make build}
	{
	  This action will build the boot device.

	  \example{\$ make build}
	}

\command{make install}
	{
	  This action will install the kaneton stuff on the boot device.

	  \example{\$ make install}
	}

\command{make export}
	{
	  This action will build a kaneton package.

	  \example{\$ make export}

	  \example{\$ make export-k5}
	}

\command{make view}
	{
	  This action will generate and display a kaneton document.

	  \example{\$ make view}

	  \example{\$ make view-develop}

	  \example{\$ make view-papers::kaneton}
	}

\command{make record}
	{
	  This action will record a real-time session.

	  \example{\$ make record}

	  \example{\$ make record-sessions::test.ts}
	}

\command{make play}
	{
	  This action will play a recorded session.

	  \example{\$ make play}

	  \example{\$ make play-sessions::prototypes.ts}
	}

\command{make info}
	{
	  This action will display general kaneton information.

	  \example{\$ make info}
	}

%%
%% licence       kaneton licence
%%
%% project       kaneton
%%
%% file          /home/buckman/kaneton/view/books/ia32-virtual/bootloader.tex
%%
%% created       matthieu bucchianeri   [sat sep  2 11:40:18 2006]
%% updated       matthieu bucchianeri   [sat sep  2 11:40:32 2006]
%%

%
% bootloader
%

\chapter{Bootloader}

%
% ---------- header -----------------------------------------------------------
%
% project       kaneton
%
% license       kaneton
%
% file          /home/enguerrand/...view/book/architecture/mips64/memory.tex
%
% created       enguerrand raymond   [tue oct  7 08:54:35 2008]
% updated          [tue apr 28 14:30:30 2009]
%

%
% memory
%

\chapter{Memory management}

This chapter describes the memory management developped for MIPS64 microprocessors.

\subsection{Bootloader}

\subsection{kaneton kernel}

%%%
%% licence       kaneton licence
%%
%% project       kaneton
%%
%% file          /home/buckman/kaneton/view/books/ia32-virtual/event.tex
%%
%% created       matthieu bucchianeri   [sat sep  2 11:40:54 2006]
%% updated       matthieu bucchianeri   [sun sep  3 12:24:43 2006]
%%

%
% event management
%

\chapter{Event management}

%
% overview of exceptions, interrupts and hardware interrupts
%

\section{Overview of exceptions, interrupts and hardware interrupts}


%%
% ---------- header -----------------------------------------------------------
%
% project       kaneton
%
% license       kaneton
%
% file          /home/buckman/kaneton/view/book/architecture/ia32/timer.tex
%
% created       matthieu bucchianeri   [tue aug 14 20:09:50 2007]
% updated       matthieu bucchianeri   [thu aug 16 17:50:18 2007]
%

%
% timer management
%

\chapter{Timer management}

This chapter describes the architecture-dependent code driving the
kaneton timer manager.

\newpage

%
% fixed frequency interrupts
%

\section{Fixed frequency interrupts}

The kaneton timer manager requires the function
\textit{timer\_handler()} to be called every
\textit{TIMER\_MS\_PER\_TICK} milliseconds.

There are three components able to throw an interrupt at fixed rate:

\begin{itemize}
  \item
    The microprocessor itself, through the local APIC;
  \item
    The Programmable Interval Timer (PIT);
  \item
    The Real Time Clock (RTC).
\end{itemize}

The first solution is the most accurate but needs to be calibrated,
which is a more complex process. The last solution is also very
accurate but the available frequencies need floating point computing.

To keep the code as simple as possible, we will use the PIT, which is
a very easily programmable component.

The PIT is connected to a fixed frequency oscillator of 1193180
Hz. The component act as a counter/decounter. We use the mode 2: a
value is loaded in the counter register and then decremented on each
clock tick. Once the counter register reaches zero, an interrupt is
signaled and the counter is reset to its initially loaded value.

The interrupt pin is connected to the IRQ line 0.

The very short function \textit{ibmpc\_timer\_init()} setup the PIT
channel 0 at mode 2 with an initial counter register computed as
follow:

$$
latch = \frac{IBMPC\_CLOCK\_TICK\_RATE}{\frac{1000}{TIMER\_MS\_PER\_TICK}}
$$

Where \textit{IBMPC\_CLOCK\_TICK\_RATE} is the fixed oscillator
frequency (1193180 Hz).

To finish, the glue code of \textit{timer\_initialize()} hooks the IRQ
0 to the function \textit{timer\_handler()}.

%%
% ---------- header -----------------------------------------------------------
%
% project       kaneton
%
% license       kaneton
%
% file          /home/buckman/kaneton/view/book/architecture/ia32/task.tex
%
% created       matthieu bucchianeri   [sat sep  2 11:41:15 2006]
% updated       matthieu bucchianeri   [sat dec 22 23:11:16 2007]
%

%
% task management
%

\chapter{Task management}

This chapter describes all the IA-32 specific mecanisms used in the
thread, task and scheduler managers.

These mecanisms mainly includes the manipulation of contexts. The
context is the whole state of the microprocessor. The machine
dependent operations consists in creating, maintaining and changing
the contexts.

\newpage

%
% reserving a task
%

\section{Reserving a task} % XXX

%
% preparing a new thread
%

\section{Preparing a new thread} % XXX

%
% scheduling and context switch
%

\section{Thread switching}

The thread's scheduling involves many mecanisms and operations on
contexts. The following sections describes the manipulations of
contexts and the context switching as implemented in kaneton for
IA-32.

%
% context saving and restoring
%

\subsection{Context saving and restoring}

Saving the context is done only on interrupt or exception. Restoring
the context is done only when resuming from such event.

The code of the low-level interrupt and exception handlers must be
generated using one of these four macro functions:

\begin{itemize}
\item \textit{IA32\_EXCEPTION\_PREHANDLER\_CODE}, for exceptions with error
code appended by the microprocessor ;
\item \textit{IA32\_EXCEPTION\_PREHANDLER\_NOCODE}, for exceptions without
any error code ;
\item \textit{IA32\_INTERRUPT\_PREHANDLER}, for interrupts ;
\item \textit{IA32\_IRQ\_PREHANDLER}, for interrupts with PIC aknowlegment ;
\item \textit{IA32\_SYSCALL\_PREHANDLER}, for system calls.
\end{itemize}

The macro functions generates the context management code, which
skeleton is:

\begin{itemize}
\item Switching to the thread's kernel stack (hardware performed) ;
\item Saving the context, \textit{IA32\_SAVE\_CONTEXT()} ;
\item Switching to the kernel interrupt stack ;
\item Calling the higher-level handler ;
\item Switching to the resuming thread's kernel stack ;
\item Restoring the context, \textit{IA32\_RESTORE\_CONTEXT()} ;
\item Returning from interrupt or exception.
\end{itemize}

The \textit{IA32\_SAVE\_CONTEXT} step makes a \textit{pusha} to save
all the general purpose registers and then pushes the data segment
register. Then, it changes the value of the page-directory base
register with kernel page-directory base (stored by the kernel in the
\textit{ia32\_interrupt\_pdbr} variable). Next it sets the kernel
segments (variable \textit{ia32\_interrupt\_ds}). To finish, it saves
the current stack pointer to \textit{ia32\_local\_jump\_stack} before
jumping to the interrupt dedicated stack
(\textit{ia32\_local\_interrupt\_stack}).

\textit{RESTORE\_CONTEXT} does the opposite: it switches back to the
stack \textit{ia32\_local\_jump\_stack}, it pops every registers
pushed on the stack and then switches address space to the one stored
in \textit{ia32\_local\_jump\_pdbr}.

The code of the generated handlers is then placed into the
\textit{.handler} ELF section, which will be mapped in every address
space.

We'll talk about the FPU and SSE contexts in the section
\textbf{Context switching}.

%
% accessing a thread's context
%

\subsection{Accessing a thread's context}

As the execution context of one thread is stored on its kernel stack,
it is impossible to access it directly since this kernel stack is only
mapped into the thread's address space.

Accessing a thread's context requires calling the \textit{as\_read}
and \textit{as\_write} operations.

The function \textit{ia32\_get\_context} retrieve the address of a
thread's kernel stack and uses \textit{as\_read} to read the stored
context. The function \textit{ia32\_set\_context} does the opposite,
except you can specify a mask of registers to update in the thread's
context. This requires to read the context and then to rewrite only
the masked registers, which is slower but easier for the programer.

%
% context switching
%

\subsection{Context switching}

The context of a thread is stored onto the kernel stack of the thread
when an exception or an interrupt occurs. When resuming, the context
at the top of the stack is used to resume execution.

Therefore, the implementation of \textit{sched\_switch} dependent code
consists only in few steps:

\begin{enumerate}
\item Setting the \textit{ia32\_local\_jump\_pdbr} variable to the
PDBR of the resuming thread ;
\item Setting the \textit{ia32\_local\_jump\_stack} variable to the
address of the kernel stack of the resuming thread ;
\item Setting up the TSS to switch stack to the thread's kernel stack
on interrupt ;
\item Copying the permission bitmap of the task into the active TSS
(see \textbf{I/O management} for details about the permission bitmap).
\end{enumerate}

The context switching of FPU and extended registers (SSE and SSE2) is
implemented using IA-32 facilities: when the TS bit of CR0 is set, an
exception is thrown each time a floating or extension instruction is
executed.

Thus, switching from one thread to another sets the TS bit to 1. If
the thread executes no special instruction, the FPU \& extended context
stay unchanged. If the thread executes such instruction, an exception
is thrown. The exception handler (\textit{ia32\_extended\_context\_switch}):

\begin{enumerate}
\item Save the FPU \& extended register set into the \textit{mmx\_context}
thread ;
\item Restore the FPU \& extended context of the currently executing thread ;
\item Clear the TS flag ;
\item Set \textit{mmx\_context} to the current thread id.
\end{enumerate}

\textit{mmx\_context} is a scheduler variable that contains the
identifier of the last thread that executed a FPU or extension
instruction.

Saving and restoring the FPU \& extended context is done using
dedicated instructions: \textit{fsave} and \textit{frstor} when no SSE
extensions are present, \textit{fxsave} and \textit{fxrstor} when SSE
is supported.

%
% address space switching
%

\subsection{Address space switching}

No additional operation is required, since loading the CR3 register
involves flushing the translation caches.

%
% stack switching switching
%

\subsection{Stack switching}

As explained above, the code executed by the kernel on
interrupt/exception runs on a global kernel stack.

XXX creation + TSS

%%
% ---------- header -----------------------------------------------------------
%
% project       kaneton
%
% license       kaneton
%
% file          /home/mycure/kaneton/view/lecture/kernels/interrupts/io.tex
%
% created       julien quintard   [wed may 16 19:17:35 2007]
% updated       julien quintard   [wed may 16 19:17:36 2007]
%

\pgfdeclareimage[interpolate=true,width=180pt]
                {io_mapped}
                {figures/io_mapped}

\pgfdeclareimage[interpolate=true,width=200pt]
                {io_port}
                {figures/io_port}

\pgfdeclareimage[interpolate=true,width=320pt]
                {io_hybrid}
                {figures/io_hybrid}

%
%
%

\begin{frame}
  \frametitle{I/O device}

  An I/O device is generally composed of a control register and a data buffer.

  \-

  Example of I/O devices: Hard disk, keyboard, Interrupt Controller, serial
  port\ldots

  \-

  To address the I/O devices, the microprocessor can use one of the 3
  architectures:

  \begin{itemize}
    \item memory-mapped I/O\\
      The hardware device's memory is mapped in RAM.
    \item I/O ports\\
    \item hybrid\\
  \end{itemize}

\end{frame}

%
%
%

\begin{frame}
  \frametitle{Memory-mapped I/O}

  \begin{center}
    \pgfuseimage{io_mapped}
  \end{center}

  {\bf Advantages}\\
  \begin{itemize}
    \item restricted instruction set
    \item No special protection mechanism is needed, since protection is
      already implemented at the memory level.
  \end{itemize}

  {\bf Disadvantages}\\
  \begin{itemize}
    \item Caching. Solution: disable caching (for a page)
  \end{itemize}

\end{frame}

%
%
%

\begin{frame}
  \frametitle{I/O ports}

  \begin{center}
    \pgfuseimage{io_port}
  \end{center}

  {\bf Advantages}\\
  \begin{itemize}
    \item Can accept a large number of devices without affecting main memory.
  \end{itemize}

  {\bf Disadvantages}\\
  \begin{itemize}
    \item Enlarge the instruction set with special read and write instructions.
  \end{itemize}

\end{frame}

%
%
%

\begin{frame}
  \frametitle{Hybrid I/O}

  \begin{center}
    \pgfuseimage{io_hybrid}
  \end{center}

\end{frame}

%
%
%

\begin{frame}
  \frametitle{I/O performing}

  I/O are asynchronous:

  \begin{itemize}
    \item in emission: a hardware device may send data at any time
    \item in reception: the kernel cannot find out when the device is ready to
      receive data
  \end{itemize}

  \-

  There are 2 methods to handle asynchronous events:

  \begin{itemize}
    \item Programmed I/O (polling)
    \item Interrupt-based I/O
  \end{itemize}

\end{frame}

%
%
%

\begin{frame}
  \frametitle{Programmed I/O}



\end{frame}


%
%
%

\begin{frame}
  \frametitle{Interrupt-based I/O}



\end{frame}

%
%
%

\begin{frame}
  \frametitle{Direct Memory Access (DMA)}

  To support DMA, the hardware must have a DMA controller. This controller
  provides the following registers:

  \begin{itemize}
    \item The address register determiness the transfert's memory address.
    \item count register specifies the transfer size.
    \item control register contains the I/O port and the transfert direction.
  \end{itemize}

  \-

  {\bf Advantages}\\
  \begin{itemize}
    \item DMA minimize the microprocessor load
    \item decrease the interrupt rate
  \end{itemize}

  {\bf Disadvantages}\\
  \begin{itemize}
    \item The DMA controller is much slower than the CPU. Also, when the CPU,
      it could assume the transfert faster.
    \item 
  \end{itemize}

\end{frame}

%%
% ---------- header -----------------------------------------------------------
%
% project       kaneton
%
% license       kaneton
%
% file          /home/buckman/kan...view/book/architecture/ia32/syscalls.tex
%
% created       matthieu bucchianeri   [tue aug 14 20:10:52 2007]
% updated       matthieu bucchianeri   [thu aug 16 19:33:03 2007]
%

%
% system calls
%

\chapter{System calls}

This chapter describes the implementation details of the low-level
system calls used in the message manager.

\newpage

%
% software interrupts and register parameter-passing
%

\section{Software interrupts and register parameter-passing}

The IA-32 microprocessors offers two facilities to implement system
calls:

\begin{itemize}
  \item
    The software interrupts with the \textit{INT} instruction;
  \item
    The \textit{SYSENTER} and \textit{SYSEXIT} instructions.
\end{itemize}

We use the first technique as it is widely used and supported on any
IA-32 microprocessor.

The fastest way to pass parameters to a system call is to use the
registers. We can pass 32 bits per register. As the message primitives
take many arguments, it is difficult to put all of them in the
registers. For this reason, we have to economize one register by
associating one interrupt per system call.

%
% kernel-side code
%

\subsection{Kernel-side code}

The function \textit{ia32\_syscalls\_init()} calls
\textit{event\_reserve()} for each system call as defined in the table
in the \textit{System calls map} section.

In every handler, we use the interrupted context pointer to get the
register values and to set the output values.

\begin{verbatim}
ia32_syshandler_register()
{
  type_1    arg1;
  type_2    arg2;
  t_status   res;

  arg1 = ia32_context->eax;
  arg2 = ia32_context->ebx;

  res = function(arg1, arg2);

  ia32_context->eax = res;
}
\end{verbatim}

%
% user-side code
%

\subsection{User-side code}

The user-side system call code put the arguments in registers, call
the corresponding software interrupt and get back the result(s) from
registers.

\begin{verbatim}
syscall(arg1, arg2)
{
  eax <- arg1
  ebx <- arg2

  int 56

  res <- eax

  return (res);
}
\end{verbatim}

%
% system calls map
%

\section{System calls map}

As explained just above, one software interrupt slot is allocated for
each system call. Here is the table of these indexes, the associated
system call and the register mapping of the parameters.

\begin{center}
\begin{tabular}{| c | c | c |}
\hline
Soft-int & Syscall           & Registers \\
\hline
56       & message\_register & \textit{eax} = type; \textit{ebx} = size \\
\hline
57       & message\_send     & \textit{eax,ebx,ecx,edx} = node; \textit{esi} = type; \textit{edi} = data; \textit{ebp} = size \\
\hline
58       & message\_transmit & \textit{eax,ebx,ecx,edx} = node; \textit{esi} = type; \textit{edi} = data; \textit{ebp} = size \\
\hline
59       & message\_throw    & XXX \\
\hline
60       & message\_receive  & \textit{eax} = type; \textit{ebx} = data; \textbf{output}: \textit{ebx} = size; \textit{ecx,edx,esi,edi} = node \\
\hline
61       & message\_grab     & XXX \\
\hline
62       & message\_poll     & \textit{eax} = type; \textit{ebx} = data; \textbf{output}: \textit{ebx} = size; \textit{ecx,edx,esi,edi} = node \\
\hline
63       & message\_state    & XXX \\
\hline
64       & message\_wait     & XXX \\
\hline
\end{tabular}
\end{center}

The register \textit{eax} is always used as the result output register
(of type \textit{t\_error}).

%
% upcalls
%

\section{Upcalls}

XXX
%%
% ---------- header -----------------------------------------------------------
%
% project       kaneton
%
% license       kaneton
%
% file          /home/mycure/kaneton/view/book/kaneton/references.tex
%
% created       julien quintard   [tue jun 19 15:36:51 2007]
% updated       julien quintard   [sun dec 16 17:33:44 2007]
%

%
% ---------- references -------------------------------------------------------
%

\chapter{References}
\label{chapter:references}

This chapter contains a listing of the kaneton documents that readers might
be interested in as well.

\newpage

%
% ---------- text -------------------------------------------------------------
%

This document aimed at presenting the kaneton microkernel and its organisation.

The books listed below go further by describing more precisely a specific
component of the kaneton microkernel:

\begin{enumerate}
  \item
    \textbf{The kaneton microkernel :: core}

    \-

    That document focuses on the design and implementation of the kaneton
    microkernel core through the different managers;
  \item
    \textbf{The kaneton microkernel :: \textit{architecture}}

    \-

    Theses documents tackle the implementation of kaneton on the given
    \textit{architecture}. Additionally, these documents contain information
    on platforms coupled with this architecture.

    \-

    Thes documents can be found in two forms, either public or private. Indeed,
    these documents contain information students must not have access to.

    \-

    The architectures currently documented are listed below:

    \begin{itemize}
      \item
	\textbf{ia32}

	\-

	\textit{Platforms}: \textbf{ibm-pc}
    \end{itemize}
\end{enumerate}

%%
%% licence       kaneton licence
%%
%% project       kaneton
%%
%% file          /home/mycure/kaneton/view/papers/kaneton/bibliography.tex
%%
%% created       julien quintard   [mon may  8 18:35:35 2006]
%% updated       julien quintard   [mon may  8 20:38:56 2006]
%%

%
% bibliograpy
%

\chapter{Bibliography}

This chapter contains the bibliography.

%
% text
%

\begin{thebibliography}{0}
  \bibitem{AST-SCO}
    \textbf{Structured Computer Organization};
    by
    \textit{Andrew S. Tanenbaum}
  \bibitem{AST-CN}
    \textbf{Computer Networks};
    by
    \textit{Andrew S. Tanenbaum}
  \bibitem{AST-OSDI}
    \textbf{Operating Systems: Design and Implementation};
    by
    \textit{Andrew S. Tanenbaum, Albert S Woodhull}
  \bibitem{AST-MOS}
    \textbf{Modern Operating Systems};
    by
    \textit{Andrew S. Tanenbaum}
  \bibitem{AST-DOS}
    \textbf{Distributed Operating Systems};
    by
    \textit{Andrew S. Tanenbaum}
  \bibitem{AST-DSPP}
    \textbf{Distributed Systems: Principles and Paradigms};
    by
    \textit{Andrew S. Tanenbaum, Maarten van Steen}
  \bibitem{NAL-DA}
    \textbf{Distributed Algorithms};
    by
    \textit{Nancy A. Lynch}
\end{thebibliography}


\end{document}
